%%%%%%%%%%%%%%%%%%%%%%%%%%%%%%%%%%%%%%%%%%%%%%%%%%%%%%%%%%%%%%%%%%%%%%%%
% Fichier .tex Principal pour le Rapport de Projet RI
% Style NAACL 2016
%%%%%%%%%%%%%%%%%%%%%%%%%%%%%%%%%%%%%%%%%%%%%%%%%%%%%%%%%%%%%%%%%%%%%%%%

\documentclass[11pt,a4paper]{article}
\usepackage{naacl2016} % Assurez-vous d'avoir ce fichier .sty
\usepackage{times}
\usepackage{latexsym}
\usepackage[utf8]{inputenc} % Encodage UTF-8 pour les caractères français
\usepackage[T1]{fontenc}    % Encodage des fontes
\usepackage[french]{babel}  % Langue française (si nécessaire pour césures etc.)
\usepackage{url}
\usepackage{graphicx}       % Pour inclure des images/graphiques
\usepackage{amsmath}        % Packages mathématiques
\usepackage{amsfonts}
\usepackage{amssymb}
\usepackage{booktabs}       % Pour de jolis tableaux (\toprule, \midrule, \bottomrule)
\usepackage{multirow}       % Pour fusionner des cellules dans les tableaux
\usepackage{natbib}         % Pour la bibliographie style auteur-date
\usepackage{microtype}      % Améliore la typographie et les coupures de mots
\usepackage{float}          % Pour mieux contrôler le placement des flottants si 

\naaclfinalcopy % À décommenter pour la version finale


% Définir le chemin vers le fichier .bib (si différent du dossier courant)
% Assurez-vous que le fichier references.bib est dans le même dossier ou spécifiez le chemin
% \bibliography{references}

\title{Rapport de synthèse de la partie Informatique} % Titre suggéré, adaptez si besoin

% Mettez vos noms et affiliations ici
\author{Dominique S. Loyer\\
  Travail présenté le 28 avril 2025 \\
  au comité, dans le cadre d'un \\
  cours de doctorat, Examen doctoral, à l'UQAM \\
  loyer.dominique@courrier.uqam.ca \\}

\date{today}
\documentclass[10pt, twocolumn]{article}
\usepackage[utf8]{inputenc}
\usepackage{fontenc}
\usepackage[french]{babel}
\usepackage{amsmath}
\usepackage{amsfonts}
\usepackage{amssymb}
\usepackage{graphicx}
\usepackage{lipsum} % For dummy text, remove in final version if not needed
\usepackage{geometry}
\geometry{a4paper, margin=1in}
\usepackage{natbib}
\bibliographystyle{apalike}
\usepackage{booktabs} % For better table rules
\usepackage{caption} % For table captions
\usepackage{adjustbox} % To adjust table size if needed

\title{Rapport Scientifique : Analyse Approfondie de l'Interprétabilité, des Systèmes de Recommandation, et Comparaison BMA/Bagging}
\author{Analyse Basée Exclusivement sur les Références Fournies} % Author line can be adjusted or removed
\date{\today}

\begin{document}

\maketitle

\section*{Résumé (Abstract)}
Ce rapport présente une analyse approfondie de trois domaines clés de l'apprentissage automatique : l'interprétabilité des modèles (XAI), les systèmes de recommandation (RecSys), et une comparaison des méthodes d'ensemble Bayesian Model Averaging (BMA) et Bagging. L'analyse se fonde \textit{exclusivement} sur un corpus de huit articles de recherche spécifiés : \citet{Doran2017}, \citet{Burkart2021}, \citet{Ko2022}, \citet{Koren2011}, \citet{Domingos2012}, \citet{Solomonoff1964}, \citet{Wasserman2000}, et \citet{Domingos1997}. L'objectif concernant l'interprétabilité (Q1) est de définir et évaluer ce concept, d'analyser ses enjeux éthiques et pratiques, et de détailler les principales méthodes (fonctionnement, forces, limites) selon \citet{Doran2017} et \citet{Burkart2021}. L'objectif relatif aux systèmes de recommandation (Q2) est de définir ces systèmes, de présenter les techniques existantes, et de comparer en détail deux d'entre elles (Filtrage Basé Contenu vs. Filtrage Collaboratif par Facteurs Latents) selon \citet{Ko2022} et \citet{Koren2011}. Enfin, l'objectif concernant la comparaison BMA/Bagging (Q3) est de comparer en profondeur ces deux approches, en examinant leurs fondements théoriques, principes de fonctionnement, et implications selon \citet{Wasserman2000}, \citet{Domingos1997}, \citet{Domingos2012}, et \citet{Solomonoff1964}. Les principaux résultats synthétisent les définitions nuancées de l'interprétabilité, les forces et faiblesses des méthodes XAI, les mécanismes et défis des systèmes de recommandation, et les fondements distincts (Bayésien vs. fréquentiste/bootstrap) ainsi que les compromis (biais/variance, incertitude) entre BMA et Bagging, tels que présentés dans la littérature analysée. Les conclusions sont intrinsèquement limitées au périmètre des huit articles fournis.

\section{Introduction}

L'omniprésence croissante des modèles d'apprentissage automatique (ML) dans des secteurs critiques tels que la santé, la finance, la sécurité et les services en ligne soulève des questions fondamentales quant à leur fonctionnement interne, leur fiabilité et leur impact.[1, 2, 3, 4] La complexité inhérente à de nombreux algorithmes modernes, notamment les réseaux de neurones profonds, conduit souvent à des modèles perçus comme des « boîtes noires » (black boxes).[1, 3, 5] Cette opacité rend leur interprétation difficile, mais elle est pourtant essentielle pour établir la confiance, garantir l'équité, valider la robustesse et assurer la conformité éthique et légale des décisions prises par ces systèmes.[3, 6, 7]

Dans ce contexte, l'interprétabilité (ou explicabilité, XAI) des modèles ML n'est plus une simple option technique mais une nécessité pressante. Comprendre \textit{comment} un modèle parvient à une prédiction et \textit{pourquoi} il prend une décision spécifique est primordial, en particulier dans les domaines sensibles où les conséquences peuvent être significatives sur le plan financier, social ou personnel.[1, 2, 3, 6, 7] Les enjeux éthiques liés à l'équité, la non-discrimination, la responsabilité et la transparence exigent des mécanismes permettant d'auditer et de justifier le comportement des algorithmes.[2, 6, 8]

Parallèlement, face à l'explosion du volume d'informations disponibles en ligne, les systèmes de recommandation (RecSys) sont devenus des outils technologiques incontournables.[9, 10, 11, 12, 13] Que ce soit pour suggérer des films, des produits, des articles d'actualité ou des offres d'emploi, ces systèmes visent à filtrer le contenu et à proposer des suggestions personnalisées, adaptées aux goûts et besoins de chaque utilisateur. La compréhension fine de leurs mécanismes sous-jacents, tels que le filtrage basé sur le contenu ou le filtrage collaboratif, est cruciale pour améliorer leur pertinence, leur diversité et leur robustesse face à des défis comme la sparsité des données ou le problème du démarrage à froid (cold start).

Enfin, pour améliorer la performance prédictive et la stabilité des modèles ML, les méthodes d'ensemble se sont largement répandues. Parmi elles, le Bayesian Model Averaging (BMA) et le Bagging (Bootstrap Aggregating) proposent des approches distinctes pour combiner les prédictions de plusieurs modèles.[4, 14, 15, 16, 17, 18, 19] Le BMA s'inscrit dans un cadre Bayésien rigoureux pour gérer l'incertitude liée au choix du modèle, tandis que le Bagging utilise le rééchantillonnage bootstrap pour réduire la variance des prédicteurs instables. Comprendre leurs fondements théoriques respectifs, leurs mécanismes opérationnels et leurs implications sur le compromis biais-variance et la gestion de l'incertitude est essentiel pour choisir l'approche la plus adaptée à un problème donné.

Ce rapport a pour objectif de répondre de manière exhaustive et détaillée à trois questions spécifiques concernant ces domaines, en s'appuyant \textit{strictement} et \textit{exclusivement} sur l'analyse approfondie d'un corpus de huit articles de recherche fondamentaux fournis : \citet{Doran2017} et \citet{Burkart2021} pour l'interprétabilité ; \citet{Ko2022} et \citet{Koren2011} pour les systèmes de recommandation ; et \citet{Wasserman2000}, \citet{Domingos1997}, \citet{Domingos2012}, et \citet{Solomonoff1964} pour la comparaison BMA/Bagging. Il synthétisera les définitions, mécanismes, avantages, inconvénients et comparaisons des concepts et méthodes abordés dans ces articles spécifiques, en respectant une structure scientifique standard : état de l'art, méthodologie d'analyse, résultats et limitations, perspectives et travaux futurs, et conclusion.

\section{État de l'art (State of the Art)}

Cette section établit les fondements conceptuels nécessaires à la compréhension des analyses ultérieures, en se basant exclusivement sur les articles fournis pour chaque question.

\subsection{Concepts Fondamentaux de l'Interprétabilité en Apprentissage Automatique (Q1)}

\subsubsection{Définition et Nécessité}
L'interprétabilité, ou explicabilité, en apprentissage automatique (souvent désignée par XAI - Explainable Artificial Intelligence) vise à rendre les processus de décision et les prédictions des modèles ML compréhensibles par les humains.[1, 2, 3, 5] Cette nécessité découle de la complexité croissante des modèles, en particulier les modèles d'apprentissage profond, qui fonctionnent souvent comme des « boîtes noires », masquant la logique interne qui relie les entrées aux sorties.[1, 3, 5]

La demande d'explicabilité est particulièrement forte dans les domaines sensibles où les décisions automatisées ont des conséquences importantes, comme le diagnostic médical, l'octroi de crédits, la justice pénale ou le recrutement.[1, 2, 3, 6, 7] Dans ces contextes, la confiance dans le système est primordiale. Les utilisateurs, qu'ils soient experts du domaine, régulateurs ou personnes affectées par la décision, ont besoin de comprendre comment une conclusion a été atteinte pour pouvoir l'évaluer, la valider et potentiellement la contester.[3, 6, 7] L'explicabilité est donc essentielle pour la validation des modèles, le débogage, l'identification et la correction des biais potentiels, la garantie de l'équité, et l'assurance de la conformité aux normes éthiques et légales, comme le droit à l'explication stipulé par le Règlement Général sur la Protection des Données (RGPD) de l'Union Européenne.[2, 3]

Selon \citet{Burkart2021}, le besoin d'explicabilité émerge souvent lorsqu'il existe un certain degré d'« incomplétude » dans la modélisation du problème.[1, 2, 3, 5] Cette incomplétude signifie que certains aspects importants du problème, tels que les notions humaines d'équité, d'éthique ou de causalité, ne peuvent pas être entièrement et formellement encodés dans la fonction objectif ou les données d'entraînement du modèle.[2, 3, 5] L'explication sert alors de pont entre la logique purement statistique ou corrélative du modèle et les exigences contextuelles et normatives du monde réel.

\subsubsection{Conceptualisations Multiples}
Les définitions et les approches de l'explicabilité varient. \citet{Doran2017} proposent une conceptualisation basée sur différentes perspectives ou capacités offertes par un système d'IA [6, 7, 8, 20] :
\begin{itemize}
    \item \textbf{Systèmes Opaques :} Ce sont des systèmes dont les mécanismes internes reliant les entrées aux sorties sont invisibles ou incompréhensibles pour l'utilisateur. Ils agissent comme des oracles fournissant des prédictions sans justification.[6, 7]
    \item \textbf{Systèmes Interprétables :} Dans ces systèmes, un utilisateur (généralement avec une expertise technique) peut analyser mathématiquement les mécanismes algorithmiques pour comprendre comment les entrées sont transformées en sorties. La transparence du modèle est requise. Des exemples incluent les modèles de régression linéaires, les arbres de décision de faible profondeur ou les machines à vecteurs de support (SVM) avec des noyaux simples.[6, 8]
    \item \textbf{Systèmes Compréhensibles :} Ces systèmes, même s'ils peuvent être opaques en interne, émettent des symboles (mots, visualisations, exemples) en plus de leur sortie principale. Ces symboles permettent à l'utilisateur de relier les propriétés des entrées à la sortie en utilisant ses propres connaissances et capacités de raisonnement implicites. La compréhensibilité est une notion graduelle, dépendant de la facilité avec laquelle l'utilisateur peut interpréter ces symboles. Les visualisations de données de haute dimension (comme t-SNE) ou les cartes d'activation dans les réseaux neuronaux convolutifs en sont des exemples.[6, 7, 8]
    \item \textbf{Systèmes Vraiment Explicables (Truly Explainable) :} \citet{Doran2017} introduisent cette notion comme un idéal à atteindre. Ces systèmes ne se contenteraient pas de fournir des éléments pour une explication, mais généreraient activement des explications structurées et compréhensibles par l'homme, potentiellement via l'intégration de mécanismes de raisonnement symbolique automatisé, sans nécessiter une étape finale d'interprétation ou de post-traitement par l'utilisateur.[6, 20]
\end{itemize}

De leur côté, \citet{Burkart2021} [3], tout en reconnaissant la diversité des termes, tendent à utiliser « interprétabilité » pour désigner la compréhension du fonctionnement global du modèle (souvent associé aux modèles intrinsèquement transparents ou « boîtes blanches ») et « explicabilité » pour les méthodes visant à fournir des justifications pour les prédictions de modèles complexes et opaques (« boîtes noires »), souvent de manière \textit{post-hoc} (après l'entraînement du modèle).[2, 3] Ils citent également \citet{Gilpin2018} pour une taxonomie plus détaillée distinguant différents types d'explications.[2, 3]

Cette divergence apparente dans les conceptualisations met en lumière une réalité importante : il n'existe pas une définition unique et universellement acceptée de l'XAI. L'approche de \citet{Doran2017} [6, 8] se concentre sur les capacités intrinsèques ou extrinsèques du système et ce que l'utilisateur peut en faire (analyse mathématique vs. compréhension de symboles). En revanche, la structure de l'étude de \citet{Burkart2021} [3], qui classe les méthodes selon des axes tels que modèle-spécifique vs. agnostique, global vs. local, et intrinsèque vs. post-hoc, met davantage l'accent sur les caractéristiques de la \textit{méthode} d'explication elle-même. Cette différence suggère que le choix d'une approche XAI dépend non seulement des besoins de l'utilisateur final (Quel niveau de détail? Quelle expertise technique?) mais aussi de la nature du modèle ML sous-jacent (Est-il intrinsèquement simple ou complexe?) et du type spécifique d'insight recherché (Une compréhension globale du comportement du modèle? Ou une justification pour une décision particulière?). Par conséquent, l'« explicabilité » n'est pas une propriété binaire mais un spectre de possibilités, et le contexte de l'application (le modèle, la tâche, l'utilisateur, les enjeux) est primordial pour déterminer quelle forme d'explication est la plus appropriée et la plus utile.

\subsubsection{Enjeux Éthiques et Pratiques}
Les motivations derrière la recherche en XAI sont profondément ancrées dans des enjeux éthiques et pratiques.[2, 3, 6, 7]
Sur le plan éthique, l'explicabilité est vue comme un prérequis pour garantir la responsabilité (accountability) des systèmes d'IA. Comment tenir un système pour responsable de ses décisions s'il est impossible de comprendre son raisonnement?[6] L'explicabilité permet d'évaluer si les décisions sont prises sur des bases rationnelles et conformes aux normes éthiques et légales, et non sur des corrélations fallacieuses ou des biais discriminatoires présents dans les données d'entraînement.[2, 6, 8] Elle est donc cruciale pour l'équité (fairness) et la prévention des discriminations.[2, 3]

Sur le plan pratique, l'explicabilité est nécessaire pour établir la confiance des utilisateurs et des parties prenantes.[3, 6, 7] Elle facilite le débogage des modèles en permettant d'identifier les causes d'erreurs ou de comportements inattendus. Elle permet aux experts du domaine de confronter la logique du modèle à leurs propres connaissances, potentiellement pour l'améliorer ou pour identifier ses limites de validité.[3] Elle est également essentielle pour la certification et le déploiement de systèmes d'IA dans des environnements réglementés.[1, 2, 3] Enfin, comprendre pourquoi un modèle fonctionne (ou échoue) peut conduire à de nouvelles découvertes scientifiques ou à une meilleure compréhension des phénomènes étudiés.[3, 5]

\subsection{Concepts Fondamentaux des Systèmes de Recommandation (Q2)}

\subsubsection{Définition et Objectif}
Un système de recommandation (RecSys) est une technologie conçue pour aider les utilisateurs à naviguer dans de vastes catalogues d'items (tels que des films, des livres, des produits musicaux, des articles de presse, des produits commerciaux, etc.) en leur suggérant les items les plus susceptibles de les intéresser.[9, 10, 11, 12] L'objectif principal est de lutter contre la surcharge d'information et d'améliorer l'expérience utilisateur en fournissant des recommandations personnalisées et pertinentes.[10, 11] Pour ce faire, ces systèmes analysent généralement les données relatives aux interactions passées des utilisateurs avec les items (par exemple, les évaluations, les achats, les clics) ainsi que, potentiellement, les caractéristiques des items et des utilisateurs.[10, 21] Ils prédisent l'intérêt potentiel d'un utilisateur pour des items qu'il n'a pas encore découverts et lui présentent une liste ordonnée de suggestions.[11]

\subsubsection{Principales Techniques}
\citet{Ko2022} [9, 10, 11] présentent une classification des principales approches de recommandation :

\begin{itemize}
    \item \textbf{Filtrage Basé sur le Contenu (Content-Based Filtering - CBF) :} Cette approche recommande des items similaires à ceux que l'utilisateur a appréciés dans le passé. La similarité est déterminée en analysant les attributs (caractéristiques, descriptions, métadonnées) des items. Le système construit un profil des préférences de l'utilisateur basé sur le contenu des items qu'il a évalués positivement et recherche ensuite des items dont le contenu correspond à ce profil.[10, 11]
    \item \textbf{Filtrage Collaboratif (Collaborative Filtering - CF) :} Cette technique, très populaire, repose sur l'idée que les utilisateurs qui ont eu des préférences similaires dans le passé auront probablement des préférences similaires à l'avenir. Elle identifie des utilisateurs ayant des goûts similaires à ceux de l'utilisateur cible (approche basée sur les utilisateurs ou "user-based") ou des items qui ont été évalués de manière similaire par de nombreux utilisateurs (approche basée sur les items ou "item-based"). Les recommandations sont ensuite générées en se basant sur les préférences de ces "voisins" ou sur la similarité entre items.[10, 11] \citet{Koren2011} [12, 21] distinguent plus finement les approches de voisinage (memory-based) des modèles à facteurs latents (model-based).
    \item \textbf{Systèmes Hybrides :} Ces systèmes combinent plusieurs techniques de recommandation (par exemple, CBF et CF, ou CF et des approches basées sur la connaissance) afin de tirer parti des forces de chaque méthode et de pallier leurs faiblesses respectives. L'objectif est d'améliorer la précision globale, la robustesse et la capacité à gérer des situations difficiles comme le démarrage à froid.[10, 11] \citet{Ko2022} mentionnent différentes stratégies d'hybridation (pondérée, par commutation, en cascade, mixte, etc.).[11]
\end{itemize}

\subsubsection{Techniques de Data Mining Associées}
La mise en œuvre de ces modèles de recommandation s'appuie sur diverses techniques d'exploration de données (data mining).[9, 10, 11] Le \textit{Text Mining} est souvent utilisé dans le CBF pour extraire des caractéristiques pertinentes des descriptions d'items ou des avis d'utilisateurs. L'algorithme des \textit{K-Nearest Neighbors (KNN)} est une base fréquente pour les approches CF basées sur le voisinage. Le \textit{Clustering} peut servir à regrouper des utilisateurs ou des items similaires. La \textit{Factorisation de Matrices} est une technique clé pour les modèles CF à facteurs latents, permettant de découvrir des dimensions cachées expliquant les interactions utilisateur-item et de gérer efficacement la sparsité des données. Enfin, les \textit{Réseaux de Neurones} et l'apprentissage profond sont de plus en plus utilisés pour modéliser des relations complexes et intégrer divers types de données (images, texte, séquences comportementales) dans les systèmes de recommandation modernes.[9, 10, 11]

\subsubsection{Focus sur le Filtrage Collaboratif et ses Avancées}
\citet{Koren2011} [12, 21] mettent en avant la prédominance et les avancées significatives du Filtrage Collaboratif (CF), en particulier des modèles à facteurs latents. Ils distinguent :
\begin{itemize}
    \item \textbf{Approches Basées sur le Voisinage (Neighborhood Methods) :} Calculent les recommandations en se basant sur les similarités entre utilisateurs (user-user) ou entre items (item-item), en utilisant les évaluations connues. Bien qu'intuitives, ces méthodes peuvent souffrir de problèmes de passage à l'échelle et de sparsité des données.[21]
    \item \textbf{Modèles à Facteurs Latents (Latent Factor Models) :} Ces modèles, dont la factorisation de matrices (souvent associée au terme SVD dans ce contexte) est l'exemple le plus emblématique, cherchent à expliquer les évaluations observées en projetant utilisateurs et items dans un espace latent de dimension réduite. Chaque utilisateur $u$ est représenté par un vecteur de facteurs $p_u$ et chaque item $i$ par un vecteur de facteurs $q_i$. L'interaction (par exemple, l'évaluation prédite $r̂_{ui}$) est souvent modélisée par le produit scalaire $q_i^T p_u$, potentiellement ajusté par des biais utilisateur et item.[21] Ces modèles sont considérés comme un premier choix en raison de leur précision et de leur capacité à gérer la sparsité.[12, 21]
\end{itemize}

L'analyse de la littérature suggère une évolution notable dans le domaine des RecSys. Alors que les premières approches reposaient davantage sur des similarités directes (CBF, CF basé voisinage) [10, 21], la tendance s'est clairement orientée vers des modèles capables de capturer des relations plus subtiles et latentes, comme la factorisation de matrices.[12, 21] La popularité croissante des modèles hybrides [10, 11] indique également une reconnaissance des limitations inhérentes à chaque approche prise isolément. De plus, les travaux de \citet{Koren2011} [21] soulignent l'importance cruciale d'aller au-delà des interactions statiques en intégrant des informations contextuelles dynamiques, telles que l'évolution temporelle des préférences (temporal dynamics) et le feedback implicite (comportements non évalués explicitement comme les clics ou les vues). Ces extensions, comme les modèles timeSVD++ [21], visent à surmonter les limitations classiques du CF (sparsité, démarrage à froid, nature statique des recommandations) et à fournir des prédictions plus précises et adaptatives. Cela implique que les systèmes de recommandation modernes et efficaces doivent non seulement modéliser les préférences passées mais aussi capturer la complexité et la dynamique du comportement des utilisateurs et des caractéristiques des items dans leur contexte d'interaction.

\subsection{Fondements Théoriques du BMA et du Bagging (Q3)}

\subsubsection{Principes Généraux de l'Apprentissage Automatique}
L'apprentissage automatique (Machine Learning - ML) est le domaine qui étudie comment les systèmes informatiques peuvent apprendre à réaliser des tâches à partir de données, sans être explicitement programmés pour celles-ci.[4, 22] Comme le souligne \citet{Domingos2012} [4, 23], le succès d'un système ML repose sur trois composantes essentielles :
\begin{itemize}
    \item \textbf{Représentation :} Le choix du langage formel pour représenter le classifieur (ou le modèle). Cela définit l'espace des hypothèses possibles que l'apprenant peut explorer.
    \item \textbf{Évaluation :} La fonction (ou critère) utilisée pour distinguer les bons modèles des mauvais (par exemple, la précision, l'erreur quadratique moyenne, la vraisemblance).
    \item \textbf{Optimisation :} La méthode employée pour rechercher le meilleur modèle dans l'espace des hypothèses, selon la fonction d'évaluation choisie (par exemple, la descente de gradient, la programmation génétique, etc.).
\end{itemize}
L'objectif ultime de l'apprentissage n'est pas de performer parfaitement sur les données d'entraînement, mais de \textit{généraliser} correctement à de nouvelles données jamais vues auparavant.[23] Un défi majeur est le \textit{surapprentissage (overfitting)}, qui se produit lorsque le modèle apprend des particularités ou du bruit présents dans les données d'entraînement, au détriment de sa capacité à généraliser.[23] La décomposition de l'erreur de généralisation en \textit{biais} (erreur systématique due à des hypothèses de modèle incorrectes) et \textit{variance} (instabilité du modèle due aux fluctuations dans les données d'entraînement) est un outil conceptuel fondamental pour analyser et comprendre les performances des apprenants.[4, 23] Les apprenants simples (comme les modèles linéaires) ont souvent un biais élevé mais une faible variance, tandis que les apprenants complexes (comme les arbres de décision profonds) peuvent avoir un faible biais mais une variance élevée.[23]

\subsubsection{Bayesian Model Averaging (BMA)}
Le Bayesian Model Averaging (BMA) est une approche issue de l'inférence Bayésienne pour traiter l'incertitude qui entoure le choix du "meilleur" modèle statistique parmi un ensemble de modèles candidats $M_1, \dots, M_k$.[14, 15, 19] Plutôt que de sélectionner un unique modèle et de baser toutes les inférences et prédictions sur celui-ci (ce qui ignore l'incertitude de sélection), le BMA combine les résultats de tous les modèles considérés.[19]

Le principe fondamental du BMA repose sur la loi des probabilités totales. Pour une quantité d'intérêt $\Delta$ (par exemple, un paramètre à estimer ou une prédiction future), sa distribution postérieure, compte tenu des données $D$, est une moyenne des distributions postérieures de $\Delta$ sous chaque modèle, pondérée par la probabilité postérieure de chaque modèle [15, 19] :
$$p(\Delta | D) = \sum_{j=1}^{k} p(\Delta | D, M_j) P(M_j | D)$$
Le poids $P(M_j | D)$ représente la probabilité que le modèle $M_j$ soit le vrai modèle générateur des données, après avoir observé ces données. Il est calculé à l'aide du théorème de Bayes [15, 19] :
$$P(M_j | D) = \frac{P(D | M_j) P(M_j)}{\sum_{r=1}^{k} P(D | M_r) P(M_r)}$$
Ici, $P(M_j)$ est la \textit{probabilité a priori} du modèle $M_j$, reflétant notre croyance initiale en sa plausibilité.[19] $P(D | M_j)$ est la \textit{vraisemblance marginale} (ou \textit{évidence}) du modèle $M_j$. Elle représente la probabilité des données observées sous le modèle $M_j$, intégrée sur toutes les valeurs possibles des paramètres $\theta_j$ de ce modèle, pondérée par la distribution a priori des paramètres $p(\theta_j | M_j)$ [15, 19] :
$$P(D | M_j) = \int p(D | \theta_j, M_j) p(\theta_j | M_j) d\theta_j$$
La vraisemblance marginale $P(D | M_j)$ joue un rôle crucial car elle pénalise naturellement les modèles plus complexes qui peuvent s'ajuster à un plus grand nombre de jeux de données possibles. Le rapport des vraisemblances marginales de deux modèles, $B_{ij} = P(D | M_i) / P(D | M_j)$, est appelé \textit{Bayes Factor} et mesure la force de l'évidence fournie par les données en faveur du modèle $M_i$ par rapport au modèle $M_j$.[14, 15, 19]

\subsubsection{Bagging (Bootstrap Aggregating)}
Le Bagging, proposé par Leo Breiman, est une méthode d'ensemble conçue pour améliorer la précision et surtout la stabilité des algorithmes d'apprentissage automatique, en particulier ceux qui sont "instables".[17] Un apprenant est dit instable si de petits changements dans les données d'entraînement peuvent entraîner des changements importants dans le modèle appris (exemples typiques : arbres de décision, réseaux de neurones).[17]

La procédure de Bagging est la suivante [17] :
\begin{enumerate}
    \item À partir d'un jeu de données d'entraînement original $D$ de taille $n$, générer $m$ nouveaux jeux de données $D_1, \dots, D_m$, appelés répliques bootstrap. Chaque $D_i$ est obtenu en tirant $n$ exemples de $D$ avec remise.
    \item Entraîner un modèle (classifieur ou régresseur) $h_i$ sur chaque jeu de données bootstrap $D_i$, en utilisant le même algorithme d'apprentissage de base.
    \item Agréger les prédictions des $m$ modèles $h_i$. Pour la classification, l'agrégation se fait généralement par vote majoritaire. Pour la régression, elle se fait par calcul de la moyenne des prédictions.
\end{enumerate}

\citet{Domingos1997} [17] explore spécifiquement pourquoi le Bagging fonctionne, en proposant un "compte rendu Bayésien". Il examine deux hypothèses principales :
\begin{enumerate}
    \item Le Bagging est une approximation de l'inférence Bayésienne optimale, c'est-à-dire du BMA. L'idée est que l'échantillonnage bootstrap simule l'échantillonnage à partir de la distribution postérieure des modèles, et l'agrégation par vote/moyenne approxime la moyenne sur cette distribution postérieure. Cette interprétation suppose un prior implicite particulier sur l'espace des modèles.
    \item Le Bagging ne fait pas qu'approximer une inférence Bayésienne existante, mais il \textit{modifie} (shifts) la distribution a priori effective. En créant des modèles composites par agrégation, il déplacerait la masse de probabilité vers des régions de l'espace des modèles plus robustes ou plus "correctes" pour le domaine, réduisant ainsi l'erreur due à un mauvais prior initial implicite de l'apprenant de base.
\end{enumerate}
L'article [17] suggère que la seconde hypothèse (modification du prior) est plus plausible, en particulier pour expliquer pourquoi le Bagging améliore des apprenants dont les hypothèses de base pourraient être mal adaptées au domaine.

\subsubsection{Probabilité Algorithmique et Rasoir d'Ockham}
La théorie de l'inférence inductive de \citet{Solomonoff1964} [24, 25, 26] offre une perspective fondamentale sur la probabilité et l'apprentissage à partir de données, basée sur la théorie de l'information algorithmique. La \textit{Probabilité Algorithmique (ALP)} assigne une probabilité a priori à toute séquence de données (ou à toute hypothèse calculable la générant) qui est inversement proportionnelle à la longueur de sa description la plus courte sur une machine de Turing universelle de référence.[24, 25] Plus précisément, la probabilité d'une hypothèse (programme) $p$ est $P(p) \approx 2^{-K(p)}$, où $K(p)$ est la complexité de Kolmogorov de $p$ (la longueur du plus court programme qui génère $p$).[24, 26]

Ce principe formalise mathématiquement le \textit{Rasoir d'Ockham} : les explications (modèles, programmes) plus simples sont exponentiellement plus probables a priori.[24, 25] L'inférence inductive de Solomonoff utilise ce prior universel dans un cadre Bayésien pour faire des prédictions. Bien qu'elle soit théoriquement optimale (complète, au sens où elle découvrira toute régularité calculable dans les données [25, 26]), elle est incalculable.[24, 25]

Le lien avec BMA et Bagging est conceptuel. Le BMA, en calculant la vraisemblance marginale $P(D|M_j)$, pénalise implicitement la complexité du modèle, s'alignant ainsi sur le principe d'Ockham.[19] Le Bagging, selon l'interprétation de \citet{Domingos1997} [17], en moyennant sur de nombreux modèles générés par bootstrap, pourrait converger vers une prédiction qui correspond à une région "centrale" ou "moyenne" de l'espace des modèles accessibles par l'apprenant de base à partir des données, ce qui pourrait, dans certains cas, correspondre à des solutions plus simples ou plus robustes. La théorie de Solomonoff fournit ainsi une justification ultime, bien qu'incalculable, à la préférence pour la simplicité qui sous-tend de nombreuses approches en ML, y compris implicitement les méthodes d'ensemble.

L'examen de ces trois approches révèle différentes manières d'aborder l'incertitude en apprentissage automatique. Le BMA, tel que décrit par \citet{Wasserman2000} [14, 19], s'attaque de front à l'\textit{incertitude du modèle} : étant donné un ensemble fini de modèles candidats, lequel est le plus soutenu par les données? Il utilise l'appareil Bayésien complet (priors sur les modèles et les paramètres, évidence marginale) pour quantifier cette incertitude via les probabilités postérieures $P(M_j|D)$. Le Bagging, comme expliqué par \citet{Domingos1997} [17], aborde principalement l'\textit{incertitude due à la variabilité de l'échantillon d'entraînement} et à l'\textit{instabilité de l'algorithme d'apprentissage}. Il utilise le bootstrap pour simuler cette variabilité et réduit la variance des prédictions en agrégeant les résultats de modèles entraînés sur ces échantillons simulés.[18] L'interprétation Bayésienne suggérée par Domingos tente de relier cette procédure algorithmique à une forme d'inférence Bayésienne approximative ou à une modification du prior effectif.[17] Enfin, la Probabilité Algorithmique de \citet{Solomonoff1964} [24, 25] traite d'une forme d'\textit{incertitude plus fondamentale} concernant la nature même du processus générateur des données, en liant la plausibilité d'une hypothèse à sa complexité algorithmique intrinsèque.

Cette distinction a des implications pratiques. Le BMA est théoriquement bien adapté lorsque l'on dispose d'un ensemble clairement défini de modèles structurellement différents et que l'on souhaite évaluer leur plausibilité relative de manière rigoureuse. Le Bagging est particulièrement indiqué lorsque l'on utilise un algorithme d'apprentissage connu pour être instable (comme les arbres de décision profonds) et que l'objectif principal est de réduire la variance et d'améliorer la robustesse des prédictions face aux variations des données d'entraînement. La théorie de Solomonoff, bien qu'incalculable, fournit un idéal théorique et une justification profonde au principe de parcimonie. Le choix entre BMA et Bagging dépend donc de la source principale d'incertitude que l'on cherche à modéliser ou à atténuer, ainsi que des hypothèses que l'on est prêt à faire (spécification de priors explicites pour BMA vs. procédure algorithmique pour Bagging).

\section{Méthodologie Analytique (Analytical Methodology)}

Cette section détaille l'analyse comparative des méthodes spécifiques pour chaque question, en se fondant sur les descriptions, forces et limitations identifiées dans les articles sources. L'objectif est d'extraire les nuances, les compromis et les relations entre les différentes approches présentées dans la littérature fournie.

\subsection{Analyse Détaillée des Méthodes d'Interprétabilité Sélectionnées (Q1)}

En se basant principalement sur l'étude de \citet{Burkart2021} [3], nous analysons quatre catégories de méthodes d'interprétabilité fréquemment citées : LIME, SHAP, l'extraction de règles, et l'importance des caractéristiques. Ces méthodes illustrent bien la diversité des approches (locales/globales, agnostiques/spécifiques, post-hoc).

\subsubsection{LIME (Local Interpretable Model-agnostic Explanations)}
\begin{itemize}
    \item \textbf{Fonctionnement :} LIME cherche à expliquer la prédiction d'un modèle boîte noire pour une instance spécifique en construisant un modèle interprétable (souvent linéaire ou un arbre de décision simple) qui approxime localement le comportement du modèle complexe autour de cette instance. Pour ce faire, il génère un voisinage de points de données synthétiques en perturbant l'instance originale, obtient les prédictions du modèle boîte noire pour ces points, et pondère ces points en fonction de leur proximité avec l'instance originale. Le modèle interprétable est ensuite entraîné sur cet ensemble de données pondéré. Les caractéristiques ou coefficients de ce modèle local servent d'explication.[3]
    \item \textbf{Forces :} Sa principale force est son caractère \textit{agnostique} : il peut être appliqué à n'importe quel modèle de classification ou de régression. Il vise une \textit{fidélité locale}, c'est-à-dire que l'explication est conçue pour être fidèle au comportement du modèle complexe dans le voisinage immédiat de l'instance expliquée. Les explications produites (souvent sous forme de poids de caractéristiques) sont généralement considérées comme \textit{intuitives}.[3]
    \item \textbf{Limites :} La validité de l'explication est strictement \textit{locale} et ne reflète pas nécessairement le comportement global du modèle. Les résultats peuvent être sensibles à la manière dont les perturbations sont générées et à la définition du voisinage (taille, métrique de proximité). Une certaine \textit{instabilité} a été observée, où des instances très similaires peuvent parfois recevoir des explications différentes. Enfin, la qualité de l'explication dépend du choix et de la capacité du modèle interprétable local.[3]
\end{itemize}

\subsubsection{SHAP (SHapley Additive exPlanations)}
\begin{itemize}
    \item \textbf{Fonctionnement :} SHAP est une approche unifiée pour l'interprétation des prédictions de modèles, basée sur les valeurs de Shapley issues de la théorie des jeux coopératifs. Elle attribue à chaque caractéristique une valeur d'importance (la valeur de Shapley) qui représente sa contribution marginale moyenne à la prédiction, en considérant toutes les combinaisons possibles (coalitions) de caractéristiques. La somme des valeurs de Shapley pour une prédiction donnée correspond à la différence entre la prédiction réelle pour cette instance et la prédiction moyenne sur l'ensemble des données. SHAP propose différentes implémentations, dont KernelSHAP (agnostique) et TreeSHAP (efficace pour les modèles basés sur les arbres).[3]
    \item \textbf{Forces :} SHAP repose sur des \textit{fondements théoriques solides} et possède des propriétés désirables (efficacité, symétrie, additivité), garantissant une attribution "équitable" de l'importance. Il permet des explications \textit{locales} (pour une instance) et \textit{globales} (en agrégeant les valeurs de Shapley sur l'ensemble des données). Il est souvent considéré comme plus \textit{cohérent} que LIME. TreeSHAP offre une \textit{efficacité computationnelle} pour les arbres et les forêts aléatoires.[3]
    \item \textbf{Limites :} Le calcul exact des valeurs de Shapley est \textit{computationnellement coûteux}, en particulier pour les modèles non-arborescents et un grand nombre de caractéristiques (KernelSHAP utilise des approximations). L'interprétation des valeurs de Shapley, bien que théoriquement fondée, peut être moins immédiate pour des non-experts. De plus, les calculs standards reposent souvent sur une \textit{hypothèse d'indépendance des caractéristiques} qui n'est pas toujours vérifiée en pratique, ce qui peut affecter la précision des attributions lorsque les caractéristiques sont corrélées.[3]
\end{itemize}

\subsubsection{Extraction de Règles}
\begin{itemize}
    \item \textbf{Fonctionnement :} Ces techniques visent à transformer un modèle complexe (souvent une boîte noire) en un ensemble de règles logiques (typiquement sous forme SI-ALORS) qui sont plus facilement compréhensibles par les humains. \citet{Burkart2021} [3] distinguent les approches \textit{pédagogiques}, qui traitent le modèle complexe comme un oracle et apprennent les règles en observant ses entrées-sorties, des approches \textit{décompositionnelles}, qui tentent d'extraire les règles en analysant la structure interne du modèle (par exemple, les poids d'un réseau de neurones ou la structure d'un arbre).
    \item \textbf{Forces :} Le principal avantage est la \textit{haute interprétabilité} potentielle des règles générées, qui peuvent fournir des explications claires et directes sur la logique de décision. Cela peut augmenter la \textit{transparence} et faciliter la \textit{validation par des experts} du domaine qui peuvent vérifier si les règles correspondent à leurs connaissances.[3]
    \item \textbf{Limites :} Le défi majeur est le \textit{compromis entre fidélité et interprétabilité}. Obtenir un ensemble de règles qui imite fidèlement un modèle complexe peut nécessiter un grand nombre de règles très spécifiques, ce qui nuit à l'interprétabilité. Inversement, un ensemble de règles simple peut ne pas capturer toute la complexité du modèle original. L'extraction peut être difficile pour des modèles très non-linéaires. Les méthodes décompositionnelles sont souvent \textit{spécifiques au modèle}, tandis que les méthodes pédagogiques dépendent de la qualité et de la représentativité de l'échantillonnage des données.[3]
\end{itemize}

\subsubsection{Importance des Caractéristiques (Feature Importance)}
\begin{itemize}
    \item \textbf{Fonctionnement :} Ces méthodes visent à quantifier l'influence ou la contribution de chaque caractéristique d'entrée sur les prédictions du modèle. L'\textit{importance globale} évalue l'influence moyenne sur l'ensemble du modèle, tandis que l'\textit{importance locale} se concentre sur une prédiction spécifique. Les techniques varient : analyse des coefficients (modèles linéaires), calcul du gain d'information ou de la réduction d'impureté (modèles basés sur les arbres), mesure de l'impact sur la performance du modèle lors de la permutation des valeurs d'une caractéristique (agnostique au modèle), ou analyse des gradients de la sortie par rapport aux entrées.[3]
    \item \textbf{Forces :} Les scores d'importance sont souvent \textit{simples à calculer et à comprendre}, fournissant un aperçu rapide des facteurs les plus influents. Certaines méthodes, comme l'importance par permutation, sont \textit{agnostiques au modèle}. L'identification des caractéristiques importantes peut aider à la \textit{compréhension générale} du phénomène modélisé et guider la \textit{sélection de caractéristiques} pour construire des modèles plus simples.[3]
    \item \textbf{Limites :} Ces méthodes indiquent généralement l'ampleur de l'influence mais \textit{pas sa direction} (un score élevé ne dit pas si la caractéristique augmente ou diminue la prédiction). Elles traitent souvent les caractéristiques de manière isolée et \textit{ignorent les interactions} complexes entre elles. L'importance d'une caractéristique peut être \textit{dépendante du contexte} des autres caractéristiques. Enfin, il y a un risque de \textit{confusion entre corrélation et causalité} : une caractéristique peut être jugée importante simplement parce qu'elle est corrélée à la vraie cause, sans être causale elle-même.[3]
\end{itemize}

\subsubsection{Évaluation de la Qualité des Explications}
\citet{Burkart2021} [3] soulignent la difficulté d'évaluer objectivement la "qualité" d'une explication, car celle-ci dépend fortement du contexte, de l'audience et des objectifs. Ils mentionnent plusieurs approches :
\begin{itemize}
    \item \textbf{Évaluations Humaines :} Idéalement, l'évaluation devrait impliquer des utilisateurs finaux réalisant des tâches réelles ou simplifiées.
    \item \textbf{Tâches Proxy et Métriques :} En l'absence d'évaluations humaines coûteuses, des métriques indirectes sont souvent utilisées. Pour les explications basées sur des règles ou des arbres, des mesures de complexité (nombre de règles, profondeur de l'arbre, nombre de caractéristiques utilisées) ou de sparsité sont courantes. La fidélité (à quel point l'explication correspond au modèle original) est un autre critère important.
    \item \textbf{Cadres d'Évaluation :} Des cadres comme PDR (Predictive accuracy, Descriptive accuracy, Relevancy) tentent de structurer l'évaluation selon plusieurs dimensions. Des checklists et des scores basés sur des jugements d'experts peuvent également être utilisés.
\end{itemize}
Le caractère subjectif et contextuel de ce qui constitue une "bonne" explication reste un défi majeur pour le domaine.[3]

L'analyse comparative de ces méthodes révèle des tensions inhérentes au domaine de l'XAI. Par exemple, LIME [3] privilégie la simplicité et l'applicabilité locale, au risque d'une représentation infidèle du comportement global du modèle et d'une certaine instabilité. SHAP [3], en revanche, offre une base théorique solide et une cohérence entre explications locales et globales, mais au prix d'une complexité de calcul potentiellement élevée et d'une interprétation moins immédiate des valeurs de Shapley. L'extraction de règles [3] vise une interprétabilité maximale via un format logique familier, mais se heurte souvent à un compromis difficile avec la fidélité au modèle original, surtout s'il est très complexe. Les méthodes d'importance des caractéristiques [3] sont souvent les plus simples à mettre en œuvre et à comprendre, mais elles fournissent une vue agrégée qui peut masquer des interactions cruciales ou des comportements locaux spécifiques. Ce panorama montre qu'il n'existe pas de méthode d'explication universellement supérieure. Le choix doit être guidé par un arbitrage conscient entre la portée souhaitée de l'explication (locale ou globale), le degré de précision requis par rapport au modèle boîte noire (fidélité), la facilité de compréhension pour l'audience cible (interprétabilité), et les contraintes computationnelles.

\subsection{Analyse Détaillée et Comparaison de Deux Techniques de Recommandation (Q2)}

Nous comparons ici en détail deux approches fondamentales des systèmes de recommandation, en nous basant sur \citet{Ko2022} et \citet{Koren2011} : le Filtrage Basé sur le Contenu (CBF) et le Filtrage Collaboratif basé sur les Facteurs Latents (Latent Factor CF). Ce choix permet de contraster une approche centrée sur les items et leurs attributs avec une approche centrée sur les interactions latentes entre utilisateurs et items.

\subsubsection{Filtrage Basé sur le Contenu (CBF)}
\begin{itemize}
    \item \textbf{Fonctionnement :} Le CBF recommande à un utilisateur des items similaires à ceux qu'il a aimés dans le passé. La similarité est calculée en comparant les attributs ou le "contenu" des items. Le système apprend un profil des préférences de l'utilisateur (par exemple, quels genres de films, quels auteurs de livres) basé sur les caractéristiques des items qu'il a évalués positivement. Ensuite, il recommande des items non encore vus dont les caractéristiques correspondent à ce profil.[11] Des techniques comme TF-IDF (Term Frequency-Inverse Document Frequency) pour l'analyse de texte, ou des classifieurs apprenant à distinguer les items aimés des non-aimés en fonction de leurs attributs, sont couramment utilisées.[10, 11]
    \item \textbf{Données Utilisées :} Nécessite principalement des informations sur le contenu des items (attributs, descriptions, mots-clés, genres, etc.) et un historique des préférences de l'utilisateur (évaluations, likes) pour construire son profil.[11]
    \item \textbf{Forces :}
        \begin{itemize}
            \item \textit{Indépendance Utilisateur :} Les recommandations pour un utilisateur ne dépendent pas des autres utilisateurs, ce qui évite le problème du démarrage à froid pour les nouveaux items (tant qu'ils ont une description).[11]
            \item \textit{Recommandations de Niche :} Peut recommander des items peu populaires ou spécifiques, du moment qu'ils correspondent au profil de l'utilisateur.
            \item \textit{Explicabilité :} Les recommandations peuvent être facilement expliquées en se référant aux attributs de l'item qui correspondent au profil de l'utilisateur (ex: "Recommandé car vous aimez les films de science-fiction réalisés par...").[11]
        \end{itemize}
    \item \textbf{Limites :}
        \begin{itemize}
            \item \textit{Sur-spécialisation :} Tend à recommander des items très similaires à ceux déjà connus, limitant la découverte de nouveaux centres d'intérêt (manque de sérendipité).[11]
            \item \textit{Dépendance du Contenu :} Nécessite des descriptions riches et structurées des items. Difficile à appliquer si le contenu est difficile à analyser (ex: musique, images sans métadonnées).
            \item \textit{Démarrage à Froid Utilisateur :} Nécessite un historique suffisant des préférences de l'utilisateur pour construire un profil fiable.[11]
            \item \textit{Déclin Relatif :} Son utilisation en tant que technique autonome a diminué au profit des approches collaboratives et hybrides.[10, 11]
        \end{itemize}
\end{itemize}

\subsubsection{Filtrage Collaboratif par Facteurs Latents (Latent Factor CF)}
\begin{itemize}
    \item \textbf{Fonctionnement :} Cette approche, popularisée notamment par la factorisation de matrices (ex: SVD), postule que les préférences des utilisateurs sont déterminées par un petit nombre de facteurs latents (non observés directement). Elle décompose la matrice (généralement très creuse) des interactions utilisateur-item (ex: évaluations) en deux matrices de facteurs de plus faible dimension : une matrice représentant les utilisateurs dans l'espace des facteurs ($P$) et une matrice représentant les items dans ce même espace ($Q$). Chaque ligne $p_u$ de $P$ est le vecteur de facteurs de l'utilisateur $u$, et chaque colonne $q_i$ de $Q$ est le vecteur de facteurs de l'item $i$. L'évaluation $r_{ui}$ qu'un utilisateur $u$ donnerait à un item $i$ est prédite par le produit scalaire de leurs vecteurs de facteurs, $r̂_{ui} \approx p_u^T q_i$. Souvent, des termes de biais (biais moyen global, biais utilisateur $b_u$, biais item $b_i$) sont ajoutés pour améliorer la précision : $r̂_{ui} = \mu + b_i + b_u + p_u^T q_i$.[21] Les vecteurs de facteurs et les biais sont appris en minimisant une fonction d'erreur (ex: erreur quadratique) sur les évaluations connues, souvent avec une régularisation pour éviter le surapprentissage, via des algorithmes comme la descente de gradient stochastique (SGD) ou les moindres carrés alternés (ALS).[21]
    \item \textbf{Données Utilisées :} Repose principalement sur la matrice des interactions utilisateur-item, typiquement les évaluations explicites. Cependant, des avancées majeures consistent à intégrer également le \textit{feedback implicite} (clics, historique d'achat, temps de visionnage) et les \textit{dynamiques temporelles} (évolution des goûts et de la popularité).[21]
    \item \textbf{Forces :}
        \begin{itemize}
            \item \textit{Précision :} Atteint généralement une très bonne précision prédictive, souvent considérée comme l'état de l'art, en particulier sur les grands jeux de données creux.[12, 21]
            \item \textit{Gestion de la Sparsité :} La factorisation de matrices gère nativement les données manquantes et est efficace même lorsque les utilisateurs n'ont évalué qu'une petite fraction des items.[21]
            \item \textit{Sérendipité :} Capable de découvrir des associations latentes et de recommander des items inattendus mais pertinents, allant au-delà des similarités évidentes de contenu.
            \item \textit{Pas besoin de contenu :} Ne nécessite pas d'informations descriptives sur les items.
        \end{itemize}
    \item \textbf{Limites :}
        \begin{itemize}
            \item \textit{Démarrage à Froid :} Difficulté à faire des recommandations pour les nouveaux utilisateurs ou les nouveaux items qui n'ont pas (ou peu) d'interactions enregistrées, car leurs vecteurs de facteurs ne peuvent pas être appris.[21]
            \item \textit{Explicabilité :} Les facteurs latents appris sont souvent difficiles à interpréter en termes sémantiques clairs, rendant les recommandations moins transparentes ("pourquoi cet item m'est-il recommandé?").[21]
            \item \textit{Dépendance des Données :} La qualité des recommandations dépend fortement de la quantité et de la qualité des données d'interaction disponibles.
        \end{itemize}
\end{itemize}

\subsubsection{Comparaison Directe}

Le tableau \ref{tab:cbf_vs_cf} synthétise la comparaison entre CBF et LF-CF.

\begin{table*}[ht]
\centering
\caption{Comparaison du Filtrage Basé Contenu (CBF) et du Filtrage Collaboratif par Facteurs Latents (LF-CF)}
\label{tab:cbf_vs_cf}
\begin{adjustbox}{width=\textwidth}
\begin{tabular}{@{}lll@{}}
\toprule
Caractéristique & Filtrage Basé Contenu (CBF) & Filtrage Collaboratif par Facteurs Latents (LF-CF) \\ \midrule
Principe de Base & Similarité des attributs/contenu des items & Interactions latentes utilisateur-item (comportement) \\
Données Requises & Attributs items + profil/préférences utilisateur & Matrice d'interactions utilisateur-item (évaluations, implicite) \\
Gestion Cold Start (Utilisateur) & Difficile (besoin de profil) & Moins difficile (si voisins existent), mais reste un défi \\
Gestion Cold Start (Item) & Facile (si contenu disponible) & Difficile (besoin d'interactions) \\
Diversité/Sérendipité & Faible (sur-spécialisation) & Élevée (découverte de goûts latents) \\
Explicabilité & Élevée (basée sur attributs) & Faible (facteurs latents non interprétables) \\
Forces Principales & Indépendance/Niche/Explicabilité & Précision/Sparsité/Sérendipité \\
Limites Principales & Sur-spécialisation/Dépendance contenu & Cold Start/Explicabilité/Dépendance interactions \\ \bottomrule
\end{tabular}
\end{adjustbox}
\footnotesize{\textit{Source :} Synthèse basée sur \citet{Ko2022} [11] et \citet{Koren2011}.[21]}
\end{table*}

L'analyse comparative met en évidence la complémentarité des deux approches. Le CBF [11] excelle lorsque le contenu est roi et que l'explication est importante, mais il peine à élargir les horizons de l'utilisateur. Le LF-CF [21], quant à lui, brille par sa capacité à modéliser des préférences complexes et à découvrir de nouveaux centres d'intérêt à partir des seuls comportements, mais il bute sur le manque de données initiales et l'opacité de ses recommandations. Cette complémentarité explique en grande partie pourquoi la tendance observée dans la littérature [10, 12, 21] est à la domination du filtrage collaboratif (surtout sous sa forme latente) pour la performance prédictive, mais souvent en combinaison avec d'autres techniques (notamment CBF ou basé connaissance) au sein de systèmes hybrides.[10, 11] Ces systèmes hybrides cherchent à exploiter le meilleur des deux mondes : la précision et la capacité de découverte du LF-CF, augmentées par la capacité du CBF à gérer le démarrage à froid des items et à fournir des explications basées sur le contenu lorsque c'est possible. Les avancées continues dans le LF-CF, telles que l'intégration du feedback implicite et des dynamiques temporelles décrites par \citet{Koren2011} [21], visent à renforcer encore sa position centrale tout en atténuant certaines de ses faiblesses.

\subsection{Analyse Comparative Approfondie du BMA et du Bagging (Q3)}

Cette section compare le Bayesian Model Averaging (BMA) et le Bagging (Bootstrap Aggregating) en se basant sur les travaux de \citet{Wasserman2000}, \citet{Domingos1997}, \citet{Domingos2012}, et la perspective théorique de \citet{Solomonoff1964}. Bien que les deux soient des méthodes d'ensemble visant à améliorer les performances par rapport à un modèle unique, elles reposent sur des fondements et des mécanismes distincts.

\subsubsection{Objectif Commun vs. Fondements Théoriques}
\begin{itemize}
    \item \textbf{Objectif Commun :} Les deux méthodes cherchent à obtenir une meilleure performance prédictive (ou d'estimation) et une plus grande robustesse en combinant les résultats de multiples modèles au lieu de se fier à un seul.
    \item \textbf{BMA (Fondement Bayésien) :} Le BMA est une application directe et cohérente des principes de l'inférence Bayésienne à l'incertitude du modèle.[14, 19] Il suppose qu'il existe un ensemble de modèles candidats plausibles et utilise le théorème de Bayes pour calculer la probabilité postérieure de chaque modèle, compte tenu des données observées et des croyances a priori sur les modèles et leurs paramètres.[19] L'inférence finale intègre explicitement cette incertitude sur le choix du modèle.
    \item \textbf{Bagging (Fondement Fréquentiste/Bootstrap) :} Le Bagging est une procédure algorithmique basée sur la technique statistique du bootstrap (rééchantillonnage avec remise).[17] Son objectif premier est de réduire la variance des estimateurs ou prédicteurs instables. L'interprétation Bayésienne proposée par \citet{Domingos1997} [17] est une tentative \textit{a posteriori} de relier cette procédure à la théorie Bayésienne, suggérant qu'elle pourrait approximer une forme de BMA ou modifier le prior effectif, mais le mécanisme de base reste ancré dans le rééchantillonnage fréquentiste.
\end{itemize}

\subsubsection{Principe de Fonctionnement}
\begin{itemize}
    \item \textbf{BMA :}
        \begin{enumerate}
            \item Définir un ensemble de modèles candidats $\mathcal{M} = \{M_1, \dots, M_k\}$.
            \item Spécifier une distribution a priori sur les modèles, $P(M_j)$, et une distribution a priori sur les paramètres $\theta_j$ de chaque modèle, $p(\theta_j | M_j)$.
            \item Calculer la vraisemblance marginale (évidence) de chaque modèle : $P(D | M_j) = \int p(D | \theta_j, M_j) p(\theta_j | M_j) d\theta_j$.[19]
            \item Calculer la probabilité postérieure de chaque modèle via le théorème de Bayes : $P(M_j | D) \propto P(D | M_j) P(M_j)$.[19]
            \item Calculer la prédiction (ou l'estimation) finale comme une moyenne des prédictions de chaque modèle, pondérée par leur probabilité postérieure : $\hat{\Delta}_{BMA} = \sum_{j=1}^{k} \hat{\Delta}_j P(M_j | D)$.[19]
        \end{enumerate}
    \item \textbf{Bagging :}
        \begin{enumerate}
            \item Générer $m$ échantillons bootstrap $D_1, \dots, D_m$ à partir des données originales $D$.[17]
            \item Entraîner un modèle $h_i$ sur chaque $D_i$ en utilisant le même algorithme d'apprentissage.[17]
            \item Agréger les prédictions des $m$ modèles $h_i$ par vote majoritaire (classification) ou moyenne (régression).[17]
        \end{enumerate}
\end{itemize}

\subsubsection{Gestion de l'Incertitude et Effet sur Biais/Variance}
\begin{itemize}
    \item \textbf{BMA :} Adresse explicitement l'\textit{incertitude sur le choix du modèle}.[14, 19] En moyennant sur l'espace des modèles considérés, il vise à produire une prédiction qui est, en moyenne, meilleure et moins risquée que celle issue d'un seul modèle sélectionné. Il ne cible pas spécifiquement la réduction de la variance au sens statistique classique, mais cherche plutôt le meilleur compromis global en intégrant l'incertitude structurelle.
    \item \textbf{Bagging :} Adresse principalement l'\textit{incertitude due à la variabilité de l'échantillon} et l'\textit{instabilité de l'apprenant}.[17, 18] Son effet principal est la \textit{réduction de la variance} des prédicteurs.[18] Pour les apprenants à faible biais et haute variance (comme les arbres de décision non élagués), le Bagging peut significativement améliorer la performance en stabilisant les prédictions. Il a généralement peu d'effet, voire un effet négatif, sur le biais [17] (bien que certains travaux plus récents, non inclus ici, nuancent ce point [16]). L'article de \citet{Domingos2012} [4, 23] fournit le cadre général de la décomposition biais-variance pour comprendre ces effets. Le Bagging ne fournit pas intrinsèquement une mesure de l'incertitude du modèle au sens Bayésien.
\end{itemize}

\subsubsection{Forces et Limites}
\begin{itemize}
    \item \textbf{BMA :}
        \begin{itemize}
            \item \textit{Forces :} Cohérence théorique Bayésienne [14], prise en compte explicite et principielle de l'incertitude du modèle [19], potentiellement optimal sous certaines conditions [19], intègre naturellement le rasoir d'Ockham via l'évidence marginale.[19]
            \item \textit{Limites :} Forte dépendance vis-à-vis des priors (sur les modèles et les paramètres) [19], coût computationnel élevé (calcul des intégrales pour l'évidence marginale, souvent intractable) [19], nécessité de définir et limiter l'espace des modèles $\mathcal{M}$ [19], difficultés avec les priors impropres.[19]
        \end{itemize}
    \item \textbf{Bagging :}
        \begin{itemize}
            \item \textit{Forces :} Simplicité conceptuelle et d'implémentation [17], très efficace empiriquement pour réduire la variance des apprenants instables [18], robustesse, ne nécessite pas la spécification de priors explicites.[17]
            \item \textit{Limites :} Moins fondé théoriquement (du point de vue Bayésien strict) [17], l'agrégation uniforme n'est pas toujours optimale (des schémas de pondération existent mais complexifient l'approche [27]), peut ne pas améliorer (voire dégrader) les performances des apprenants stables (biais élevé).[17]
        \end{itemize}
\end{itemize}

\subsubsection{Synthèse Comparative}

Le tableau \ref{tab:bma_vs_bagging} résume les points clés de la comparaison.

\begin{table*}[ht]
\centering
\caption{Analyse Comparative du BMA et du Bagging}
\label{tab:bma_vs_bagging}
\begin{adjustbox}{width=\textwidth}
\begin{tabular}{@{}lll@{}}
\toprule
Caractéristique & Bayesian Model Averaging (BMA) & Bagging (Bootstrap Aggregating) \\ \midrule
Fondement Théorique & Bayésien (Théorème de Bayes) & Fréquentiste (Bootstrap) / Algorithmique \\
Mécanisme Principal & Pondération par probabilités postérieures $P(M_j|D)$ & Rééchantillonnage + Agrégation (vote/moyenne) \\
Gestion de l'Incertitude & Incertitude du modèle (parmi $\mathcal{M}$) & Variabilité échantillon / Instabilité apprenant \\
Effet Principal sur Erreur & Compromis biais-variance via moyenne & Réduction de la variance (pour apprenants instables) \\
Complexité Computationnelle & Élevée (calcul évidences marginales) & Modérée (entraînement de $m$ modèles) \\
Sensibilité (ex: Priors) & Élevée (priors modèles \& paramètres) & Faible (pas de priors explicites) \\
Forces Clés & Cohérence théorique, gestion incertitude modèle & Simplicité, robustesse, réduction variance \\
Limites Clés & Coût calcul, sensibilité priors, espace $\mathcal{M}$ & Moins fondé, pas optimal pour stables, agrégation uniforme \\ \bottomrule
\end{tabular}
\end{adjustbox}
\footnotesize{\textit{Source :} Synthèse basée sur \citet{Wasserman2000} [14, 15, 19], \citet{Domingos1997} [17, 18], et \citet{Domingos2012}.[23]}
\end{table*}

La comparaison entre BMA et Bagging illustre une dualité fondamentale dans les approches de l'apprentissage automatique et de la gestion de l'incertitude. Le BMA incarne l'approche Bayésienne "orthodoxe" [19] : on définit un espace d'hypothèses (les modèles $M_j$), on spécifie nos croyances a priori (les $P(M_j)$ et $p(\theta_j|M_j)$), on met à jour ces croyances avec les données via le calcul des probabilités postérieures $P(M_j|D)$, et on prend une décision (prédiction) qui tient compte de cette distribution de croyances mise à jour. C'est une approche basée sur la modélisation probabiliste explicite de l'incertitude. Le Bagging, en revanche, représente une approche plus algorithmique et fréquentiste.[17] Face à un problème (la variabilité des prédictions d'un apprenant instable), on propose une procédure (rééchantillonner et agréger) qui fonctionne bien en pratique. La justification théorique, qu'elle soit basée sur la réduction de variance [18] ou sur une interprétation Bayésienne \textit{a posteriori} [17], vient souvent après coup pour expliquer le succès de l'algorithme. Les défis computationnels et la sensibilité aux priors du BMA [19] contrastent avec la relative simplicité d'implémentation du Bagging.[17]

En conclusion, le choix entre BMA et Bagging n'est pas seulement technique, il reflète aussi une philosophie. Si l'objectif est une modélisation rigoureuse de l'incertitude structurelle entre plusieurs modèles bien définis, et si les contraintes computationnelles et la spécification des priors sont gérables, le BMA offre un cadre cohérent. Si l'objectif principal est de stabiliser un algorithme d'apprentissage spécifique connu pour sa variance élevée et d'obtenir une amélioration robuste et simple des performances, le Bagging est souvent un choix pragmatique et efficace. La tentative de \citet{Domingos1997} de fournir un "compte rendu Bayésien" du Bagging peut être vue comme une tentative de jeter un pont entre ces deux mondes, suggérant que même une procédure algorithmique simple comme le Bagging pourrait implicitement capturer certains des bénéfices d'une approche Bayésienne plus formelle, potentiellement en explorant et en moyennant efficacement sur une partie pertinente de l'espace des modèles.

\section{Résultats et Limitations (Results and Limitations)}

Cette section synthétise les principaux résultats issus de l'analyse comparative des articles fournis pour chaque question et souligne les limitations inhérentes à cette étude, qui se base exclusivement sur le corpus documentaire spécifié.

\subsection{Synthèse des Résultats}

\subsubsection{Q1 : Interprétabilité des Modèles d'Apprentissage Automatique}
L'analyse des travaux de \citet{Doran2017} et \citet{Burkart2021} révèle que l'interprétabilité (ou explicabilité, XAI) est un concept crucial mais multidimensionnel et sans définition unique consensuelle. \citet{Doran2017} distinguent les systèmes selon leurs capacités : opaques, interprétables (analyse mathématique possible), et compréhensibles (émission de symboles interprétables par l'utilisateur), tout en visant des systèmes "vraiment explicables" intégrant un raisonnement automatisé.[6] \citet{Burkart2021} tendent à différencier interprétabilité (compréhension globale du modèle) et explicabilité (justification post-hoc pour boîtes noires).[3] Cette divergence conceptuelle montre que l'approche XAI dépend du contexte (modèle, utilisateur, objectif).

La nécessité de l'XAI est motivée par des enjeux éthiques (équité, responsabilité, non-discrimination) et pratiques (confiance, débogage, validation, conformité légale) pressants, surtout dans les domaines sensibles.[3, 6]

Les méthodes d'interprétabilité principales (LIME, SHAP, extraction de règles, importance des caractéristiques) présentent des compromis clairs.[3] LIME offre une explication locale agnostique mais potentiellement instable. SHAP fournit une base théorique solide pour l'attribution locale et globale mais peut être coûteux. L'extraction de règles vise une haute interprétabilité mais lutte avec le compromis fidélité/complexité. L'importance des caractéristiques est simple mais souvent superficielle, ignorant les interactions. Le choix d'une méthode implique donc un arbitrage entre portée (locale/globale), fidélité au modèle original, et facilité de compréhension [Insight 4]. L'évaluation de la qualité des explications reste un défi majeur, souvent subjective et dépendante du contexte.[3]

\subsubsection{Q2 : Systèmes de Recommandation}
Les systèmes de recommandation (RecSys) visent à atténuer la surcharge d'information en suggérant des items pertinents aux utilisateurs.[11] L'étude de \citet{Ko2022} et \citet{Koren2011} montre que les principales techniques incluent le Filtrage Basé sur le Contenu (CBF), le Filtrage Collaboratif (CF), et les Systèmes Hybrides.[11]

Le CBF se base sur la similarité des attributs des items aimés par l'utilisateur.[11] Il est explicable et gère bien le démarrage à froid des items, mais souffre de sur-spécialisation et nécessite des descriptions de contenu riches. Le CF, en particulier l'approche par Facteurs Latents (LF-CF) mise en avant par \citet{Koren2011} [21], modélise les interactions utilisateur-item via des dimensions latentes apprises à partir des évaluations (ou feedback implicite). Le LF-CF atteint une haute précision et favorise la sérendipité, mais peine avec le démarrage à froid (utilisateurs et items) et offre peu d'explicabilité directe.

La comparaison [Insight 5] révèle une tendance vers la domination du LF-CF pour sa performance, souvent complété par des approches hybrides pour pallier ses faiblesses, notamment le démarrage à froid. Les avancées récentes se concentrent sur l'intégration de données contextuelles comme le temps et le feedback implicite pour améliorer encore la pertinence et l'adaptabilité des recommandations.[21]

\subsubsection{Q3 : Comparaison BMA vs. Bagging}
Le BMA et le Bagging sont deux méthodes d'ensemble visant à améliorer la performance des modèles ML, mais elles diffèrent fondamentalement dans leurs principes et objectifs [Insight 6].

Le BMA, décrit par \citet{Wasserman2000}, est une approche purement Bayésienne qui gère l'\textit{incertitude du modèle} en calculant la probabilité postérieure de chaque modèle candidat et en moyennant leurs prédictions pondérées par ces probabilités.[19] Il repose sur le théorème de Bayes, les priors, et le calcul de l'évidence marginale. Il est théoriquement fondé mais computationnellement coûteux et sensible aux choix des priors.[19]

Le Bagging, analysé par \citet{Domingos1997} sous un angle Bayésien, est une procédure algorithmique basée sur le bootstrap qui vise principalement à \textit{réduire la variance} des apprenants instables en agrégeant les prédictions de modèles entraînés sur des répliques bootstrap des données.[17, 18] Il est simple à implémenter et souvent efficace empiriquement, mais sa justification théorique est moins directe que celle du BMA. L'interprétation Bayésienne suggère qu'il pourrait approximer BMA ou modifier le prior effectif.[17]

Le choix entre les deux dépend de la source d'incertitude principale à adresser : l'incertitude structurelle entre modèles (BMA) ou l'instabilité de l'apprenant due à l'échantillonnage (Bagging). La théorie de \citet{Solomonoff1964} sur la probabilité algorithmique [25] fournit un cadre ultime reliant probabilité et complexité (Rasoir d'Ockham), pertinent conceptuellement pour les deux approches qui, à leur manière, favorisent ou aboutissent à des solutions d'une certaine forme de "simplicité" ou de "robustesse".

\subsection{Limitations de l'Analyse}

Il est essentiel de reconnaître les limitations de ce rapport, qui découlent directement des contraintes de l'exercice :
\begin{itemize}
    \item \textbf{Périmètre Exclusif des Sources :} L'intégralité de l'analyse, des définitions et des conclusions repose \textit{uniquement} sur les informations extraites des huit articles de recherche spécifiés. Aucune connaissance externe, aucun développement ultérieur dans les domaines concernés, ni aucune autre perspective théorique ou méthodologique n'ont été pris en compte. Les conclusions sont donc valides uniquement dans le cadre de ce corpus limité.
    \item \textbf{Profondeur et Actualité Variables :} Les articles fournis couvrent différentes périodes et abordent les sujets avec des niveaux de détail variables. Par exemple, l'article de \citet{Domingos1997} sur le Bagging [17] pose des questions sur son interprétation Bayésienne mais ne fournit pas une analyse exhaustive ou une preuve définitive.[17] De même, certaines références datent de plusieurs années (ex: \citet{Wasserman2000}, \citet{Koren2011}), et des avancées significatives ont pu avoir lieu depuis dans ces domaines.
    \item \textbf{Absence de Validation Empirique Nouvelle :} Ce rapport est une synthèse analytique et critique de la littérature fournie. Il ne contient aucune nouvelle expérience, simulation ou validation empirique des méthodes discutées. Les affirmations sur les performances relatives (ex: précision de LF-CF, réduction de variance par Bagging) sont basées sur les déclarations et résultats rapportés dans les articles sources.
    \item \textbf{Interprétation Analytique :} Bien que l'analyse s'efforce de rester fidèle aux textes, l'identification des relations entre les concepts, la synthèse des arguments et la formulation des "insights" (présentés ici comme des liens logiques ou des implications) impliquent nécessairement une part d'interprétation et de structuration par l'auteur de ce rapport (agissant en tant que persona expert).
\end{itemize}

Ces limitations signifient que ce rapport doit être considéré comme une photographie détaillée de la manière dont ces sujets spécifiques sont abordés dans ce corpus documentaire précis, plutôt qu'un état de l'art exhaustif et à jour des domaines concernés.

\section{Perspective et Travaux Futurs (Perspective and Future Work)}

L'analyse approfondie des huit articles de référence, bien que limitée à ce corpus, met en lumière des tendances importantes et suggère plusieurs directions pour les recherches futures dans les domaines de l'interprétabilité, des systèmes de recommandation et des méthodes d'ensemble.

\subsection{Implications des Résultats et Tendances}

L'impératif croissant d'interprétabilité [3, 6] pousse la recherche soit vers la conception de modèles intrinsèquement plus simples et transparents (quitte à sacrifier potentiellement une partie de la performance brute), soit vers le développement de méthodes d'explication \textit{post-hoc} plus fiables, robustes et dont la qualité peut être évaluée de manière plus rigoureuse.[3] Le défi consiste à fournir des explications qui soient non seulement fidèles au modèle mais aussi réellement utiles et compréhensibles pour différentes audiences.

Dans le domaine des systèmes de recommandation, l'évolution vers des modèles capables de capturer des interactions latentes complexes (comme LF-CF) [12, 21] et d'intégrer des informations contextuelles dynamiques (temps, feedback implicite) [21] semble se poursuivre. Cependant, les défis persistants du démarrage à froid et de l'explicabilité [21] indiquent un besoin continu de développer des approches hybrides [10] ou de nouvelles techniques qui équilibrent mieux précision, robustesse et transparence.

La comparaison entre BMA et Bagging [Insight 6] souligne un débat méthodologique fondamental en apprentissage statistique : faut-il privilégier une modélisation probabiliste rigoureuse de l'incertitude (approche Bayésienne type BMA) ou des procédures algorithmiques pragmatiques dont l'efficacité est validée empiriquement (approche type Bagging)? La recherche de ponts théoriques entre ces deux approches, comme l'interprétation Bayésienne du Bagging [17], reste un domaine d'intérêt.

\subsection{Pistes de Recherche Issues des Limitations Analysées}

Les limitations identifiées dans les articles sources et dans l'analyse comparative suggèrent plusieurs pistes pour des travaux futurs :

\subsubsection{Interprétabilité (XAI)}
\begin{itemize}
    \item \textbf{Évaluation Rigoureuse :} Développer des métriques et des protocoles standardisés pour évaluer objectivement la qualité, la fidélité, la robustesse et l'utilité des différentes méthodes d'explication, en allant au-delà des métriques proxy [3] et en intégrant davantage d'évaluations humaines dans des contextes réalistes.
    \item \textbf{Stabilité et Fidélité :} Mener des études comparatives approfondies sur la stabilité des explications fournies par des méthodes comme LIME et SHAP face à de légères perturbations des données ou du modèle, et mieux caractériser le compromis entre fidélité locale et globale.[3]
    \item \textbf{Raisonnement Intégré :} Poursuivre la recherche sur les systèmes "vraiment explicables" qui intègrent des capacités de raisonnement symbolique pour générer des justifications causales ou logiques, comme envisagé par \citet{Doran2017}.[6]
    \item \textbf{Explications Causales :} Développer des méthodes XAI qui tentent d'aller au-delà des corrélations pour fournir des explications basées sur des inférences causales, un aspect souvent désiré par les utilisateurs mais difficile à obtenir des modèles ML standards.[3]
\end{itemize}

\subsubsection{Systèmes de Recommandation (RecSys)}
\begin{itemize}
    \item \textbf{Gestion Améliorée du Cold Start :} Développer de nouvelles techniques pour les modèles à facteurs latents afin de mieux gérer le démarrage à froid pour les nouveaux utilisateurs et items, potentiellement en intégrant plus efficacement des informations externes ou des techniques hybrides.[21]
    \item \textbf{Explicabilité des Recommandations :} Concevoir des modèles de recommandation (en particulier basés sur CF) qui soient intrinsèquement plus explicables ou développer des méthodes d'explication post-hoc spécifiquement adaptées aux RecSys, sans compromettre de manière significative la précision.[21]
    \item \textbf{Intégration Contextuelle Avancée :} Explorer des manières plus sophistiquées d'intégrer divers types de contexte (social, spatial, émotionnel, etc.) dans les modèles de recommandation.[11, 21]
    \item \textbf{Dynamique à Long Terme :} Modéliser l'évolution des préférences des utilisateurs sur le long terme, en tenant compte des changements de goûts, de l'influence sociale et des événements de vie.
    \item \textbf{Équité et Biais :} Étudier et atténuer les biais potentiels (popularité, position, etc.) dans les algorithmes de recommandation et promouvoir l'équité et la diversité dans les suggestions.
\end{itemize}

\subsubsection{Méthodes d'Ensemble (BMA et Bagging)}
\begin{itemize}
    \item \textbf{Efficacité Computationnelle du BMA :} Développer des algorithmes d'approximation plus rapides et plus précis pour le calcul des évidences marginales dans le BMA, ainsi que des stratégies de recherche efficaces dans de grands espaces de modèles.[19]
    \item \textbf{Sensibilité aux Priors du BMA :} Étudier plus avant l'impact des choix de priors (modèles et paramètres) sur les résultats du BMA et développer des méthodes pour spécifier des priors objectifs ou robustes, notamment en présence de priors impropres.[19]
    \item \textbf{Théorie du Bagging :} Approfondir la compréhension théorique du Bagging, notamment en clarifiant son lien avec l'inférence Bayésienne [17] et en analysant plus finement son effet sur le biais dans différents contextes.[16]
    \item \textbf{Comparaisons Systématiques :} Réaliser des études empiriques comparant systématiquement BMA et Bagging (et d'autres méthodes d'ensemble comme le Boosting) sur une large gamme de jeux de données, de types d'apprenants et de métriques de performance, afin de mieux délimiter leurs domaines d'applicabilité respectifs.
    \item \textbf{Lien avec la Complexité Algorithmique :} Explorer comment les concepts de la probabilité algorithmique et de la complexité de Kolmogorov [25] pourraient être utilisés, même de manière approximative, pour guider la sélection de modèles ou la pondération dans les approches d'ensemble comme BMA ou pour mieux comprendre l'effet de régularisation implicite du Bagging.
\end{itemize}

Ces pistes de recherche, suggérées par l'analyse de ce corpus spécifique, montrent que malgré les avancées significatives, ces trois domaines de l'apprentissage automatique restent dynamiques et offrent de nombreuses opportunités pour des contributions théoriques et pratiques futures.

\section{Conclusion}

Ce rapport a fourni une analyse détaillée et comparative de trois sujets centraux en apprentissage automatique – l'interprétabilité des modèles, les systèmes de recommandation, et les méthodes d'ensemble BMA et Bagging – en se basant exclusivement sur un corpus restreint de huit articles de recherche fondamentaux.

Concernant l'\textbf{interprétabilité (Q1)}, l'analyse a montré qu'il s'agit d'un concept multidimensionnel essentiel pour la confiance, l'éthique et la validation des systèmes d'IA, en particulier dans les domaines sensibles.[3, 6] Les conceptualisations varient [3, 6], mais un consensus émerge sur la nécessité de dépasser les modèles "boîtes noires". Les méthodes existantes comme LIME, SHAP, l'extraction de règles et l'importance des caractéristiques offrent des solutions avec des compromis distincts en termes de portée (locale/globale), de fidélité, d'interprétabilité et de coût computationnel.[3] L'évaluation rigoureuse de la qualité des explications demeure un défi ouvert.[3]

Pour les \textbf{systèmes de recommandation (Q2)}, le rapport a défini leur rôle crucial dans la gestion de la surcharge d'information [11] et a présenté les principales approches : Filtrage Basé Contenu (CBF), Filtrage Collaboratif (CF), et Hybrides.[11] La comparaison détaillée entre CBF [11] et Filtrage Collaboratif par Facteurs Latents (LF-CF) [21] a mis en évidence leurs forces et faiblesses complémentaires. Le LF-CF domine souvent en termes de précision et de capacité de découverte, mais souffre de problèmes de démarrage à froid et d'explicabilité, justifiant la popularité des approches hybrides et l'importance d'intégrer des informations contextuelles et dynamiques.[21]

Enfin, la \textbf{comparaison entre BMA et Bagging (Q3)} a révélé deux philosophies distinctes pour l'amélioration des modèles via l'ensemble [Insight 6]. Le BMA [19], ancré dans la théorie Bayésienne, modélise explicitement l'incertitude du modèle mais fait face à des défis computationnels et de spécification de priors.[19] Le Bagging [17], basé sur le bootstrap, est une procédure algorithmique simple et efficace pour réduire la variance des apprenants instables [18], dont la connexion à l'inférence Bayésienne est un sujet d'interprétation.[17] Le choix entre les deux dépend largement de la nature de l'incertitude à traiter et des contraintes pratiques.

En conclusion, l'analyse de cette littérature spécifique, bien que limitée par son périmètre, souligne la complexité, les défis persistants et les compromis inhérents à la conception et à l'application de modèles d'apprentissage automatique qui soient à la fois performants, interprétables, pertinents et robustes. Les directions de recherche futures suggérées visent à adresser ces défis pour continuer à faire progresser le domaine de manière responsable et efficace.


% --- j) Bibliographie ---
%\section*{Références} % Utiliser section* pour ne pas numéroter la bibliographie

\nocite{*}
\bibliographystyle{acl} % Style bibliographique NAACL/ACL
\bibliography{ete25Info} % 
%\bibliography{referencesVF} % Nom de votre fichier .bib
\clearpage
\tableofcontents

\end{document}